\documentclass[11pt,a4paper]{article}
\usepackage[T1]{fontenc}
\usepackage[utf8]{inputenc}
\usepackage{lmodern}
\usepackage[a4paper,margin=25mm,headheight=14pt]{geometry}
\usepackage{amsmath,amssymb,amsthm,mathtools}
\usepackage{booktabs,tabularx,array,longtable}
\usepackage{microtype}
\usepackage{xcolor}
\usepackage{enumitem}
\usepackage{fancyhdr}
\usepackage{fancyvrb}
\usepackage{xurl}
\usepackage[colorlinks=true,linkcolor=blue!45!black,urlcolor=blue!45!black,citecolor=blue!45!black]{hyperref}
\hypersetup{pdftitle={Smaller Diophantine Equations for ZFC Consistency},pdfauthor={ChatGPT, prepared for Vladimir Reshetnikov},pdfsubject={Explicit polynomial reductions, exact certificates, and proof dependencies}}
\newtheorem{theorem}{Theorem}[section]
\newtheorem{lemma}[theorem]{Lemma}
\newtheorem{proposition}[theorem]{Proposition}
\newtheorem{corollary}[theorem]{Corollary}
\theoremstyle{definition}
\newtheorem{definition}[theorem]{Definition}
\newtheorem{remark}[theorem]{Remark}
\DeclareMathOperator{\Con}{Con}
\setlist{itemsep=3pt,topsep=5pt}
\setlength{\emergencystretch}{2em}
\allowdisplaybreaks[1]
\pagestyle{fancy}
\fancyhf{}
\fancyhead[L]{\small Diophantine equations for ZFC consistency}
\fancyhead[R]{\small Constructive reductions}
\fancyfoot[C]{\thepage}
\title{\textbf{Smaller Diophantine Equations\\for ZFC Consistency}\\[5mm]
\large Four-root separation, normalized radix positions,\\and direct certificates with recovered domains}
\author{Prepared for Vladimir Reshetnikov and the ProveIt repository\\
\small AI-assisted research by ChatGPT; proofs and exact artifacts supplied for review}
\date{23 October 2026\\\small Comparison snapshot: \texttt{7e533c33819ee1f703b69f2778a767cdc069940d}}

\begin{document}
\maketitle
\begin{abstract}
We give seven explicit ordinary Diophantine polynomials, each with a
two-way existence reduction to a fixed syntactic refutation predicate for
ZFC. The principal unrestricted-integer proposal uses fourteen unknowns,
has exact degree 540, and admits a 160-operation arithmetic circuit. A
147-operation companion has degree 568. Three positive-domain equations
retain nine unknowns while improving their immediate predecessors in
degree and circuit length; the shortest has 420 operations. A separate
direct proof-calculus construction gives a quartic in 495 nonnegative
unknowns. Recovering most coordinate signs from its own equations reduces
the unrestricted-integer lift to 977 unknowns, with 24,632 explicitly
listed monomials and coefficient height 83,869,250.

The new arguments are ordinary mathematical proofs composed with pinned
ProveIt calibration theorems. Exact finite checks reconstruct coefficients,
replay transformations, establish degree certificates, and reject
corrupted data. We supply neither a new end-to-end proof-assistant theorem
nor an object-language ZFC proof file. No global minimality or historical
priority is claimed. The contribution is a reproducible improvement of
several distinct measures within the inspected repository, together with
structural ideas that can be developed further.
\end{abstract}

\clearpage
\tableofcontents
\clearpage

\section{The result and what ``simplest'' measures}
\label{sec:overview}

An ordinary Diophantine equation is a polynomial equation with fixed
integer coefficients. Its unknowns range over a specified integer domain.
Every exponent in such a polynomial is a fixed nonnegative integer; a
variable exponent, factorial, or binomial symbol cannot remain as a
primitive in the final equation. All seven delivered artifacts meet this
requirement. Registers in the compact formulas denote polynomial
abbreviations. They are not extra existential unknowns.

There are several conflicting measures of simplicity: the number of
unknowns, total degree, arithmetic-circuit length, coefficient magnitude,
and the size of the correctness argument. A circuit can be very short
while referring to large fixed coefficients. An expanded quartic can have
modest coefficients while needing many variables. The same distinction
underlies the study of universal variable--degree pairs
\cite{bayer}, although the present problem concerns a single fixed
consistency instance rather than a universal family.

\begin{table}[ht]
\centering
\small
\setlength{\tabcolsep}{4pt}
\begin{tabular}{llrrl}
\toprule
Equation & Domain & Unknowns & Exact degree & Further measure\\
\midrule
Boundary, simple & $\mathbb Z_{>0}$ & 9 & 1,472,168 & 420 operations\\
Boundary, complete & $\mathbb Z_{>0}$ & 9 & 997,810 & 1,057 operations\\
Compact, simple & $\mathbb Z_{>0}$ & 9 & 1,626,248 & 428 operations\\
$P_{540}$ & $\mathbb Z$ & 14 & 540 & 160 operations\\
$P_{568}$ & $\mathbb Z$ & 14 & 568 & 147 operations\\
$P_{495}$ & $\mathbb Z_{\ge0}$ & 495 & 4 & 15,377 monomials\\
$P_{977}$ & $\mathbb Z$ & 977 & 4 & 24,632 monomials\\
\bottomrule
\end{tabular}
\caption{The seven delivered polynomials. The domain is part of each claim.}
\label{tab:results}
\end{table}

For unrestricted integer variables with moderate degree and a short
formula, our main proposal is $P_{540}$. Its lower-operation companion
$P_{568}$ provides a useful alternative. If positive integers are the
intended domain, the 420-operation equation is the shortest new evaluator
in this packet. The compact 428-operation form keeps a smaller fixed
coefficient tuple. If a direct proof-calculus source and an entirely
expanded polynomial with moderate coefficients are the priority, use
$P_{977}$.

The comparison is deliberately specific. At the inspected commit~\cite{proveit},
incoming fourteen-integer forms included degree 652 with 230 operations,
degree 696 with 150 operations, and degree 792 with 137 operations. Thus
$(568,147)$ improves $(696,150)$ in both measurements, and $(540,160)$
offers a still lower degree. The degree-792 form retains a smaller circuit.
Other incoming constructions use fewer than fourteen unrestricted
integers, at much higher degree. The new fourteen-variable equations are
not variable-count records.

The direct comparison is particularly large: the synchronized parent has
a quartic lift in 2,088 unrestricted integers. Our final lift has 977,
a reduction of 1,111 variables, or approximately $53.2\%$. Its coefficient
height is also slightly smaller, and the quartic degree is unchanged.
These numbers compare actual collected equations under the same domain,
not a natural-domain equation against an integer-domain one.

\subsection{The precise logical claim}

Fix the repository's sealed recursive presentation of the ZFC axioms and
its finite first-order derivation predicate. Write
\[
 \Con(\mathrm{ZFC})\equiv
 \neg\exists p\in\mathbb N\;
       \operatorname{Refutation}_{\mathrm{ZFC}}(p).
\]
This notation concerns syntactic consistency. It is not an assumption
that every theorem of ZFC is true in an intended universe.

\begin{theorem}[Ordinary consistency equivalence]
\label{thm:main}
For each delivered polynomial $P$ and its domain $D^n$ in
Table~\ref{tab:results}, the new reductions and the named inherited
calibration theorems give an ordinary proof formalizable in ZFC of
\[
 \exists\mathbf x\in D^n\;P(\mathbf x)=0
 \quad\longleftrightarrow\quad
 \exists p\in\mathbb N\;
       \operatorname{Refutation}_{\mathrm{ZFC}}(p).
\]
Consequently,
\[
 \boxed{\mathrm{ZFC}\vdash
  \bigl(\nexists\mathbf x\in D^n:\ P(\mathbf x)=0\bigr)
        \longleftrightarrow\Con(\mathrm{ZFC}).}
\]
\end{theorem}

The proof is compositional: Sections~\ref{sec:crt14},
\ref{sec:radix-origin}, \ref{sec:projected-direct}, and
\ref{sec:frame-domain} prove the new existence equivalences;
Section~\ref{sec:dependencies} identifies the imported endpoints. None of
these implications assumes consistency. In particular, the theorem does
not assert the absence of solutions. A proof of that further assertion
inside ZFC would also prove its consistency; the present biconditional
does not supply such a proof.

\subsection{Three ideas responsible for the improvements}

\paragraph{Separate the guard before combining radicals.}
The fourteen-integer equation keeps one square-root witness explicit.
Only four Pell radicals remain in the norm. Their magnitudes already
admit a strict superincreasing order under a weak scale hypothesis.
That observation removes a large auxiliary amplifier. Rational extraction
can therefore precede the CRT argument that restores the stronger cubic
scale. Centering the norm around polynomial root approximations gives
the degree-540 form.

\paragraph{Use the units digit of the coefficient layout.}
The positive nine-variable interface does not require an empty radix
position zero. A uniform origin translation preserves every source
coefficient identity, lowers the convolution height, and permits a new
exact population-count calibration. Shared powers then shorten the
literal evaluator. This is a reduction in a valid encoding, not merely a
change in the formatting of a polynomial.

\paragraph{Recover domains from the arithmetic certificate.}
The direct quartic already forces most of its coordinates to be
nonnegative. The common Pell frame supplies additional signs using only
ordered-ring inequalities, before invoking its stronger power semantics.
Only 170 coordinates need explicit nonnegativity representations.
Applying signed eliminations afterward yields the 977-integer equation.
This order prevents the usual independent representation of every
natural coordinate from dominating the variable count.

\section{Four Pell roots: fourteen integers at degrees 540 and 568}
\label{sec:crt14}

This section gives two ordinary polynomials whose fourteen coordinates all
range over the unrestricted integers. The first has exact degree $540$ and
$160$ live binary arithmetic operations; the second has degree $568$ and
$147$ operations. The fixed coefficient tuple has $615\,883$ raw minimal
unsigned bytes, with the original six integer payloads retained. Operations
count shared binary addition, subtraction, and multiplication, including
multiplication by a fixed constant. A register is a polynomial abbreviation,
not an additional quantified coordinate.

\subsection{Literal definition}

The variables, in circuit order, are
\[
 j,C,z,h,r,v,w,T,c,n,t,p_2,p_3,p_1\in\mathbb Z.
\]
Use the six fixed integers of the retained source-34/source-71 tuple:
\[
 \beta=2^{1064921}-1,\quad K_0,M_C,M_F,\quad
 A_{\mathrm c}=2^{1113},\quad B_{\mathrm c}=2^{666104}.
\]
Here $M_F$ is already shifted by $\beta$. The exact integer bytes, finite
machine-and-tile recipe, original licenses, and ordinary calibration texts
are included in the component's \texttt{upstream/} directory. This tuple is
larger than some other native tuples in ProveIt; the present determinant
scales are proved for this whole tuple and are not mixed with another one.

Define
\begin{align*}
 J&=6j+1,&q&=\beta J+1,&W&=12w+4,\\
 X&=qW,&Y&=qv,&E&=XY,\\
 a&=E+Y,&H&=4a+3,&\Delta&=a^2+H.
\end{align*}
Then define the transport and rank registers
\begin{align*}
 F&=(K_0+W)C-z(q-1),\\
 R&=(q(q-F)-C)(q^2-1)+(M_C+qM_F)J.
\end{align*}
Continue with
\begin{align*}
 k&=R+1+hE,&\eta&=c-Yk,&\zeta&=k-\eta,\\
 i&=2r+1,&I&=ic^2,&D_f&=1+\Delta I^2,\\
 Q&=\Delta^2I^2,&\delta&=Q-1,&V&=D_f(cT-R)-c,\\
 C_y&=QV^2-1,&\ell&=XY^2,&L&=\ell k,\\
 D_\tau&=1+\ell(\ell+1)k^2,&D_m&=1+\Delta c^2,\\
 b_0&=X+ac,&F_0&=(a+2)I.
\end{align*}
Put
\begin{align*}
 \operatorname{Pos}(u_1,u_2,u_3)&=1+u_1^2+u_2^2+u_3(u_3+1),\\
 G&=(A_{\mathrm c}F)(B_{\mathrm c}C)
       -(q-A_{\mathrm c}F-B_{\mathrm c}C)^2,\\
 B&=v^2G\eta\zeta-(t^2-v^2G-\eta\zeta+1)^2
       -\operatorname{Pos}(p_1,p_2,p_3),\qquad S=B^2.
\end{align*}
Let
\[
 \mathcal F_{16}(x;A_1,A_2,A_3,A_4)
 =\prod_{\varepsilon\in\{-1,1\}^4}
       \left(x-\sum_{j=1}^4\varepsilon_j\sqrt{A_j}\right).
\]
Sign invariance makes this an integer polynomial in its displayed
arguments. Its exact $495$ nonzero coefficients are verified against the
independent product of the sixteen factors.

For the centered equation define
\begin{align*}
 M&=4q^2,&U&=MH,\\
 Z&=U(n+2V+2L+k+2F_0)+2Mb_0+2H(v-q^2),\\
 N&=\delta^8\mathcal F_{16}
   \left(Z;4U^2D_f,4U^2D_\tau,4M^2D_m,\frac{4U^2C_y}{\delta}\right).
\end{align*}
\[\boxed{P_{540}=N+S+NS=0.}\]
The denominator is cleared symbolically; the literal evaluator contains
only integer polynomial operations. For the shorter companion use instead
\begin{align*}
 M&=2q^2,&U&=MH,\\
 Z&=U(n+V)+Mb_0+H(v-q^2),\\
 N&=\delta^8\mathcal F_{16}
 \left(Z;U^2D_f,U^2D_\tau,M^2D_m,\frac{U^2C_y}{\delta}\right).
\end{align*}
\[\boxed{P_{568}=N+S+NS=0.}\]
All other registers remain unchanged.

\subsection{The two guard facts}

For integers $x,y,t$ define
\[
 \operatorname{Pair}(x,y,t)=xy-(t^2-x-y+1)^2.
\]
Then
\[
 (\exists t\in\mathbb Z)\;\operatorname{Pair}(x,y,t)>0
 \quad\Longleftrightarrow\quad x>0\ \text{and}\ y>0.
\]
The forward direction follows because positive Pair forces $xy>0$, whereas
for $x=-u,y=-v$ with $u,v>0$ its subtracted square is at least
$(u+v+1)^2>uv$. For the reverse direction set $K=x+y-1$ and choose a
nearest integer square $t^2$. Between consecutive squares $m^2,(m+1)^2$
the error is at most $m$, so $(t^2-K)^2<K\le xy$; square endpoints,
including $K=1$, have zero error.

The polynomial Pos represents precisely the positive integers. Positivity
is immediate; surjectivity follows from the three-square theorem applied
to $4n-3$, with one odd and two even square coordinates. Consequently,
$B=0$ forces $v\ne0$, $G>0$, and $\eta\zeta>0$. Conversely these three
conditions admit private witnesses making $B=0$.

Since $S=B^2\ge0$, the equation
\[
 P=(1+S)(1+N)-1=0
\]
forces $S=N=0$: $1+S$ is a positive integer divisor of $1$. No sign
hypothesis about the norm $N$ is used.

\subsection{Rational extraction before the cubic scale is known}

The front end gives $q\equiv2\pmod6$ and $|q|\ge16$. Also
$W=4(3w+1)\ne0$, hence $x=|X|\ge64$. When $v=0$ we have $a=0$,
$H=3$, $\Delta=3$; when $v\ne0$, $m=|Y|\ge16$ and $|a|\ge1008$.
Thus $H\ne0$ and $\Delta\ge3$ on the entire integer domain.
The three integer radicands $D_m,D_f,D_\tau$ are positive, the last by
$\ell(\ell+1)\ge0$.

If $c=0$, then $Q=0$, $\delta=-1$, and $C_y/\delta=1$. If $c\ne0$,
then $Q\ge9$, $D_f>|c|+1$, and the integer $V=D_f(cT-R)-c$ has
$|V|\ge1$. Therefore
\[
 \frac{C_y}{\delta}=V^2+\frac{V^2-1}{Q-1}\ge1.
\]
In particular $\delta$ never vanishes, including all exceptional branches.

At a polynomial zero, the guard supplies $v\ne0$, $G>0$, and
$\eta\zeta>0$. Write $K=|k|$, $C_0=|c|$. Then $K\ge2$ and
$0<\eta/k<1$, so
\[
 K(m-1)<C_0<K(m+1),\qquad C_0>30.
\]
Let $d,f,\tau,y$ be the positive square roots of
$D_m,D_f,D_\tau,C_y/\delta$, respectively. They satisfy
\[
 \frac{d}{|H|}<\frac{C_0}{2},\qquad
 \tau>2C_0,\qquad f>2\tau.
\]
For the first inequality use $H^2-4\Delta=12a^2+8a-3$.
For the second use
$K(xm^2-1)<\tau<K(xm^2+1)$.
For the third,
\[
 d>((x-1)m-3)K(m-1)>\frac34\tau,
 \qquad f>\frac{C_0}{2}d>15d.
\]
The middle comparison follows by expanding
$m^2(x-4)-4m(x+2)+9>0$ for $x\ge64,m\ge16$.
Also $f>3C_0^2$ and $Q>f^2$.

If $cT-R=0$, then $C_0<y<C_0+1$. Otherwise
$y>|V|\ge f^2-C_0>2f$. Hence either
\[
 f,\tau,y,d/|H|\quad\text{or}\quad y,f,\tau,d/|H|
\]
is a strictly superincreasing descending list: every entry exceeds the
sum of all later entries. For example,
$y+d/|H|<1.5C_0+1<2C_0<\tau$ in the first case.
The deleted source amplifier $4H+1$ is unnecessary because the guard root
is now a separate explicit coordinate, leaving these four roots to separate.

A rational signed sum of positive roots with rational squares and strictly
superincreasing magnitudes makes every root rational. Indeed, a nontrivial
Galois automorphism changing at least one radical sign would yield a
nonempty vanishing signed subsum after subtraction, contradicting the
largest-term inequality. This argument allows dependent square classes.
At $N=0$ one conjugate factor vanishes; division by $2U$ in the centered
case, or $U$ in the companion, gives exactly such a rational sum.
Thus $d,f,\tau$ are integers.

The auxiliary rational $y$ is also integral. For any integers $Q>1,V$,
\[
 QV^2-(Q-1)y^2=1,\quad y\in\mathbb Q
 \quad\Longrightarrow\quad y\in\mathbb Z.
\]
If $y=a/b$ is reduced, then $b^2\mid Q-1$. The transformation
\[
 V'=(2Q-1)V-2(Q-1)y,\qquad
 y'=(2Q-1)y-2QV
\]
preserves the equation and denominator, keeps $V'$ integral, and for
$V>1$ gives $-V<V'<V$. Here $y>V$ and $(Q-1)y<QV$ give the two
strict inequalities. The equation excludes $V'=0$. Descent reaches
$|V|=1$, where $|y|=1$, proving the claim.

\subsection{The congruences and the consistency equivalence}

Absorb the conjugate signs into integer roots. In the centered case their
relation is
\[
 Z=2Md+2U(f+\tau+y).
\]
Since $H\equiv3\pmod q$, $q\equiv2\pmod3$, and $H$ is odd,
$\gcd(H,4q^2)=1$. Reduction modulo $M=4q^2$ yields
$2q^2\mid v-q^2$, and reduction modulo $H$ yields $H\mid d-b_0$.
Therefore
\[
 v=q^2(2v_{\rm old}+1),\quad
 Y=(2v_{\rm old}+1)q^3,\quad
 \gamma=(d-b_0)/H\in\mathbb Z.
\]
The companion with $M=2q^2$ gives the same conclusions without the
factors $2$. No cubic-scale promise was used in the preceding extraction.

These roots, the restored scale, and the two strict guards give the retained
parent's normalized signed kernel. Its original strong variables can be
restored explicitly:
\begin{align*}
 I_{\rm old}&=\Delta i c^2,& r_{\rm old}&=(\Delta i-1)/2,\\
 o_{\rm old}&=f(cT-R),&j_{\rm old}&=D_fT-1-R\Delta i^2c^3.
\end{align*}
The quotient is integral because $\Delta i$ is odd. Substitution gives
$I_{\rm old}^2=\Delta(f^2-1)=Q$ and
$fo_{\rm old}-c=cj_{\rm old}-R=V$. Private old guard witnesses are
regenerated from $G>0,\eta\zeta>0$.

Conversely, take a normalized source witness, set
$v=q^2(2v_{\rm old}+1)$, and choose the private guard witnesses by Pair
and Pos. For the centered equation choose
\[
 n=2(\gamma+f+\tau+y-V-L-F_0)-k
        -\frac{v-q^2}{2q^2}.
\]
For the companion choose
\[
 n=\gamma+f+\tau+y-V-\frac{v-q^2}{2q^2}.
\]
Both are integers and make the selected conjugate factor vanish. Every
quantity is chosen before $n$, so the construction is not circular.

We have proved an exact two-way existence reduction to the normalized
kernel. Composing with the parent's explicitly identified ordinary native
compiler and K215/refutation calibration yields, for $d=540,568$,
\[
 \mathrm{ZFC}\vdash
 \left(\neg\exists\boldsymbol x\in\mathbb Z^{14}
              \ P_d(\boldsymbol x)=0\right)
 \longleftrightarrow\operatorname{Con}(\mathrm{ZFC}).
\]
This imports ordinary mathematical theorems and their stated canonical
completeness. It does not promote finite arithmetic checks into an unbounded
semantic proof, or supply the separately open encoded-ZFC proof artifact.

\subsection{Exact degrees, including cancellation}

Treat the fixed coefficients as degree zero. For leading parts set
\begin{align*}
 q_0&=6\beta j,&W_0&=12w,&X_0&=q_0W_0,&Y_0&=q_0v,\\
 a_0&=X_0Y_0,&i_0&=2r,&R_0&=q_0^4z-q_0^3W_0C,\\
 k_0&=R_0+ha_0,&\eta_0&=-Y_0k_0,&Q_0&=a_0^4i_0^2c^4,\\
 V_0&=-a_0^2i_0^2c^4R_0,&L_0&=X_0Y_0^2k_0,&F_{00}&=a_0i_0c^2.
\end{align*}
The degrees of $\delta,U,V,L,F_0,\eta$ are $22,6,19,11,7,7$.
The guard has top part $-\eta_0^4$, so $S$ has degree $56$ and
leading part $\eta_0^8$.

For the centered equation use approximate roots
\[
 r_y=2UV,\quad r_\tau=U(2L+k),\quad
 r_f=2UF_0,\quad r_d=2Mb_0,
 \quad s=Un+2H(v-q^2).
\]
The center is their sum plus $s$. The four exact square defects are
\begin{align*}
 4U^2C_y/\delta-r_y^2&=4U^2(V^2-1)/\delta,\\
 4U^2D_\tau-r_\tau^2&=U^2(4-k^2),\\
 4U^2D_f-r_f^2&=4U^2(1-I^2),\\
 4M^2D_m-r_d^2&=4M^2(1+Hc^2-2acX-X^2).
\end{align*}
The approximate root degrees are $25,17,13,7$, and the corresponding
formal square-root correction degrees are at most $3,5,5,4$.
To justify this calculation, scale every variable by an indeterminate $u$
and work over the Laurent-series field of the rational function field in
the original coordinates. Each radicand is an approximate-root square
times $1$ plus negative-degree terms, so the binomial series supplies its
formal root with that leading sign. The delivered polynomial uses no series.

Eight conjugate factors have degree $25$, four degree $17$, two degree
$13$, and the last two degree $7$. For the final two, the leading terms
are $U_0n$ and $U_0(n+c)$, where $U_0=16q_0^2a_0$.
The product therefore has degree $308$, and the clearing factor
$\delta^8$ adds $176$. Thus $\deg N=484$, and
\[
 (P_{540})_{540}
 =2^{28}U_0^{16}Q_0^8V_0^8L_0^4F_{00}^2\eta_0^8 n(n+c).
\]
In particular the actual fixed-coefficient slice satisfies
\[
 [j^{240}w^{82}v^{94}h^{12}r^{34}c^{68}z^8n^2]P_{540}
   =2^{530}3^{322}\beta^{240}>0.
\]
The selected exponents sum to $540$, proving exact degree.

For the simpler companion only the auxiliary root is centered. Eight
factors have degree $25$ and eight have degree $17$; the norm has degree
$8\cdot25+8\cdot17+8\cdot22=512$. Now $U_0=8q_0^2a_0$, and
\[
 (P_{568})_{568}=2^8U_0^{16}Q_0^8V_0^8L_0^8\eta_0^8.
\]
Its selected coefficient is
\[
 [j^{256}w^{88}v^{104}h^{16}r^{32}c^{64}z^8]P_{568}
 =2^{520}3^{344}\beta^{256}>0.
\]
The exact audits check these entire leading forms, all four square-defect
identities, all $495$ coefficients of the generic norm, all six fixed
integer payloads against their finite reconstruction, and both literal
DAGs. Finite modular comparisons are additionally retained as transcription
diagnostics; the unbounded logical argument is the one above.

\section{Removing unused radix positions in nine-variable equations}
\label{sec:radix-origin}

This section gives three complete polynomial constructions in nine strictly
positive integers.  The improvement has two parts: an exact translation of
the finite coefficient layout, and a smaller evaluator for the translated
polynomial.  The translation is an existence-preserving change of the
encoding; the evaluator change is an identity of polynomials.  These two
claims should be kept distinct.

\begin{center}
\begin{tabular}{lrrr}
\hline
Construction & Variables & Exact degree & Binary operations\\
\hline
Original31, simple majorant &9&1,472,168&420\\
Original31, complete radical majorant &9&997,810&1,057\\
Compact54, simple majorant &9&1,626,248&428\\
\hline
\end{tabular}
\end{center}

The respective incoming comparisons are the saturated H1868 simple form
(1,474,688; 449 operations), its complete-radical form (999,518; 1,086), and
the compact54 saturated form (1,628,768; 434).  Thus the reductions relative
to those immediate predecessors are respectively 2,520 degrees and 29
operations, 1,708 degrees and 29 operations, and 2,520 degrees and six
operations.  The 420-operation form also has lower degree and fewer
operations than the 434-operation compact form, but it has a larger
coefficient payload.  These are improvements in stated measures within the
inspected repository.  They do not establish global minimality or reduce
the number of variables below nine.

The first two predecessors are in
\texttt{ConZFC\_Saturated\_Certificates\_2026-10-24.zip}; the third is in
\texttt{con-zfc-saturation-and-projection-2026-10-23.zip}.
The repository snapshot is
\texttt{7e533c33819ee1f703b69f2778a767cdc069940d}.
The included source ledger identifies these exact archives by SHA-256.

\subsection{A translation lemma for sparse coefficient layouts}

Let $\rho_1,\ldots,\rho_m$ be quadratic polynomials in $\nu$ variables,
and suppose fixed coefficient polynomials $K_0(T),K_1(T),K_2(T)$, distinct
nonnegative weights $w_i$, and distinct nonnegative target positions $q_j$
satisfy
\begin{equation}
 [T^{q_j}]\sum_{d=0}^{2}K_d(T)U_{\mathbf p}(T)^d
       =2\rho_j(\mathbf p),\qquad
 U_{\mathbf p}(T)=\sum_{i=1}^{\nu}p_iT^{w_i}.
 \label{eq:origin-extraction}
\end{equation}
The equality is a polynomial identity in all source coordinates, including
coefficients of monomials absent from the desired residual.  It is stronger
than a numerical identity at selected source tuples.

\begin{lemma}[Translation of the layout]
\label{lem:radix-origin}
Let $r,s$ be nonnegative integers such that $w_i\ge r$, $q_j\ge s$, and
every occupied exponent $e$ of $K_d$ satisfies $e+dr-s\ge0$.
Set
\[
 w_i'=w_i-r,\quad q_j'=q_j-s,\quad
 K_d'(T)=\sum_e [T^e]K_d(T)\,T^{e+dr-s}.
\]
Then these are ordinary polynomial data and satisfy the same extraction
identities with primes on $w,q,K,U$.
If $W=\max_iw_i$ and
$H=\max_d(\deg K_d+dW)$, their new values are
$W'=W-r$ and $H'=H-s$.
\end{lemma}

\begin{proof}
As a temporary identity in Laurent polynomials,
$U_{\mathbf p}'=T^{-r}U_{\mathbf p}$ and
$K_d'=T^{dr-s}K_d$.  Therefore
\[
 \sum_d K_d'(U_{\mathbf p}')^d
       =T^{-s}\sum_dK_dU_{\mathbf p}^d.
\]
Coefficient extraction proves the claim.  The hypotheses ensure that every
exponent in the actual delivered data is nonnegative, so the construction
contains no division by an unknown.  Finally
$\deg K_d'+dW'=\deg K_d+dW-s$, proving the height identity.
\end{proof}

Apply the lemma with $(r,s)=(1,2)$.  Both inherited layouts have minimum
coordinate weight one, minimum occupied linear-kernel exponent one, and
minimum occupied quadratic-kernel exponent zero.  Their constant kernels
begin above exponent two.  Consequently the transformation reads
\begin{equation}
 w_i'=w_i-1,\qquad q_j'=q_j-2,\qquad
 K_2'=K_2,\quad K_1'=K_1/T,\quad K_0'=K_0/T^2.
 \label{eq:origin-shift}
\end{equation}
The last two expressions denote removal of initial zero coefficients;
the JSON tables and the final circuits use nonnegative exponents only.
No source coordinate, source equation, or numerical coefficient changes.

There is a useful restricted optimality statement.  Among translations of
this form with $r,s\ge0$ and nonnegative resulting exponents, $r\le1$
because the minimum original weight is one, and $s\le2r$ because the
quadratic kernel has a nonzero constant coefficient.  Thus $(1,2)$ is
maximal in both parameters for these fixed layouts.  This says nothing
about changing the weights nonuniformly, changing the source residuals,
or using a different encoding architecture.

The translated coordinate weights, in source-variable order, are
\[
\begin{split}
 &(0,6,60,35,15,84,19,66,42,10,33,12,95,79,28,49,\\
 &\hspace{12mm}1,108,109,111,115,90,52,54,26,100,91,8,106,36,59).
\end{split}
\]
For the original31 source, the target positions in row order are
\[
\begin{split}
 &(1139,26,1435,1067,1834,1352,1812,612,265,1406,1401,\\
 &\hspace{12mm}887,335,1585,411,671,1137,770,1129,1016,1053,83).
\end{split}
\]
The source is the inherited positive system named \texttt{gapN86}, with
31 coordinates and 22 quadratic residuals.  Its old coordinate maps are
$P_{\rm old}=P+\eta+\zeta$ and $N_{\rm old}=N+1$; its recorded elementary
row operations are retained unchanged.  The compact source is the
separate \texttt{gapNq85} system.  Its tables and whole compact54 coefficient
tuple are kept separately; substituting individual coefficient ports
between the two systems is not justified.

\subsection{Why a coordinate may occupy position zero}

The choice $w_i>0$ was not needed by the positive-domain bit-mask argument.
For distinct weights in $\{0,\ldots,W\}$, put
\[
 \ell=\sum_iB^{w_i},\qquad N_0=B^{W+1},\qquad
 M_0=N_0-1-(b-1)\ell.
\]
Assume $b$ and $B$ are powers of two, with $b\mid B$ and $b\le B$.
The radix-$B$ digits of $M_0$ are $B-b$ at selected positions and $B-1$
elsewhere.  For $0\le g<N_0$, binary disjointness $g\mathbin{\&}M_0=0$
is therefore equivalent to
\begin{equation}
 g=\sum_i z_iB^{w_i},\qquad 0\le z_i<b.
 \label{eq:origin-mask}
\end{equation}
This digit argument includes position zero without any modification: the
allowed bits there are precisely its bottom $\log_2b$ bits.  The recovered
positive source coordinates are $p_i=z_i+1$.

It is essential to check the next arithmetic interface rather than merely
observe that the digit argument works.  The pinned Lean source
\texttt{Diophantine/Paper1982/NineRatio.lean}, retained in the packet,
states \texttt{Jones1982.Nine.ratio\_sound} and
\texttt{ratio\_complete} with the numerical hypotheses
\[
 R,N\ge8,\quad 0<b\le R,N,
\]
and with the appropriate positive carrier, margin, central-binomial
and power-of-two conditions.  Neither theorem assumes $b\mid g$ or
$b\mid R+1$.  Hence the possible nonzero units digit introduced by
\eqref{eq:origin-mask} does not violate an inherited premise.  This point
is specific to the positive nine-variable interface; the signed
Sun/Lucas construction has additional residue normalizations and is not
changed here.

The all-one source tuple is not a common zero: its fifth residual is one,
independently of the four coefficient ports.  Thus every genuine source
zero has at least one $p_i>1$, and its packed $g$ is positive even with a
weight zero.  For all sufficiently large binary $b$, its digits satisfy
\eqref{eq:origin-mask} and $g<bB^W$.  The cofinal completeness property
survives the translation.

\subsection{The fully recalibrated original31 prefix}

For this paragraph all layout data are the translated original31 data.
Write $K_d$ for their three kernels, and put
\[
 \kappa=2134939,\qquad B=2^\kappa b^2,\qquad
 T_*=\sum_{j=1}^{22}B^{q_j},\qquad U_* =\ell+g.
\]
Define
\begin{align}
 G_*&=(K_2(B)U_*+K_1(B))U_*+K_0(B),\notag\\
 N_0&=B^{116},& M_0&=N_0-1-(b-1)\ell,\notag\\
 S_*&=g+N_0(2G_*+T_*),&
 T_c&=M_0+N_0(B-2)T_*,\notag\\
 L&=2bB^{1983},&
 R&=L(S_*+T_c+1)-S_*-1,\notag\\
 N&=2^{16012032}B^{1053},&t&=bB^{115}.
 \label{eq:origin-prefix}
\end{align}
Every exponent here is fixed.  The kernel degrees are $(1406,1726,1636)$,
so $W=115$, $H=1866$, and the largest target is $1834=H-32$.
The leading kernel coefficients remain $-6,2,1$.

The inherited coefficient bound is $M=2^{2134927}$, and
$B=4096Mb^2$.  The exact artifact checker also verifies this bound on
all actual fixed coefficients.  For positive $b,g$, one has
$U_*\ge B^{115}$ and $K_2(B)>B^{1636}/2$.
Relative to $B^{1636}U_*^2$, the possible negative linear and constant
contributions are bounded by $2M/B^{25}$ and $2M/B^{460}$.
Thus $G_*>0$ unconditionally.  The low mask lies in $[0,N_0)$, so
$S_*,T_c\ge0$ and
\[
 R=(L-1)(S_*+T_c+1)+T_c\ge L-1\ge8.
\]
Also $b\le R,N$ and $N\ge8$.

On the soundness range $g<4t$, let $Z_*=B^{1867}$.
The same geometric-series estimate as in the inherited sparse-carry proof
now gives
\[
 U_*<6bB^{115},\qquad
 \frac{2G_*}{Z_*}<\frac{86Mb^2}{B-1}<\frac1{32}.
\]
The two target-mask bounds are unchanged because their headroom $H-Q=32$
is unchanged.  Together with $g<N_0$ and $M_0<N_0$, they give
\[
 S_*+T_c<N_0Z_*=B^{1983}<L.
\]
Each source convolution coefficient has magnitude at most
$993Mb^2<B/4$.  The retained signed-digit carry lemma therefore applies
with exactly the same strict inequalities.  It identifies
$S_*\mathbin{\&}T_c=0$ with a positive common source zero.

For completeness, the threshold must be recalculated.  Write
$b=2^k$, $B=2^v$, where $v=\kappa+2k$, and $L=2^u$.
There are 116 low radix digits, 31 selected low positions, and 22 high
mask digits.  Consequently
\[
 \operatorname{pop}(T_c)=138v-31k-22,\qquad
 u=1+k+1983v.
\]
The two-block population identity gives
\[
 \operatorname{pop}(R)=u+\operatorname{pop}(T_c)-\tau,
 \quad
 \tau=\operatorname{pop}(S_*)+\operatorname{pop}(T_c)
       -\operatorname{pop}(S_*+T_c)\ge0,
\]
with $\tau=0$ exactly for binary disjointness.  Now
\begin{align*}
 u+\operatorname{pop}(T_c)
   &=2121\kappa+4212k-21\\
   &=2\bigl(16012032+1053(\kappa+2k)\bigr)
     =2\log_2N,
\end{align*}
since $15\kappa-21=32024064$.
Legendre's formula therefore yields
\begin{equation}
 N^2\mid\binom{2R}{R}
 \quad\Longleftrightarrow\quad S_*\mathbin{\&}T_c=0.
 \label{eq:origin-calibration}
\end{equation}
The argument includes $b=1$.  At that value $M_0=N_0-1$ and
$0<g<N_0$, so disjointness fails.  There is no exceptional accepted
assignment at the endpoint.

For the compact54 variant replace the source tables and the whole native
coefficient tuple, and use
\[
 \kappa=864026,\quad H=2080,\quad Q=2048,\quad
 L=bB^{2197},\quad N=2^{6480184}B^{1160},\quad t=bB^{115}.
\]
The kernel degrees are $(1625,1940,1850)$, the dominance gaps are 25 and
455, and the coefficient bound is $2^{864014}$.  Its population calculation
is $u+\operatorname{pop}(T_c)=2335\kappa+4640k-22=2\log_2N$.
Thus the same proof applies.  The exact source audit checks all 22
identities separately for each of these two genuinely different sources.

\subsection{The actual nine-variable polynomials}

The unknowns are
\[
       b,g,h,s,w,C,i,j,n\in\mathbb Z_{>0}.
\]
All other names in this subsection are abbreviations, not quantified
variables.  Use either certified prefix above and define
\begin{align*}
 U&=N^2w,&Y&=N^2s,\\
 M&=RY,&P&=2M^2U,\\
 K&=R+1+h(P-1),&A&=M(U+1),\\
 B_1&=2R+1,&W_1&=bw,\\
 J&=4A-5,&Q&=3W_1C-2(W_1^2-1),\\
 m&=K^2(4t-g)-16t(C-KY)^2,&\Omega&=Jm,\\
 E_0&=J(n+W_1^2)+Q,&V_0&=PK,\\
 a_0&=(P^2-1)K^2+1,&D&=(A^2-1)C^2+1,\\
 r_0&=iJC^2,&F&=(A^2-1)Dr_0^2+1,\\
 \Lambda&=CD+1-A,&G&=A+F\Lambda,\\
 Z&=B_1+(j-1)C,&&\\
 H_1&=C+FZ,&I&=(G^2-1)H_1^2+1.
\end{align*}
For either simple construction set
\[
        a_1=DFI,\qquad V_1=ADr_0GH_1.
\]
For the complete-radical construction set
\[
 \sigma=1024C^5i,\qquad a_1=\sigma^2DFI,\qquad
 V_1=\Pi(A,C,i,Z),
\]
where $\Pi$ is the fully specified inherited 333-term integer polynomial
in \path{data/saturated_corrected_P.json}.
Its exact identity certificate is retained; it is not an unspecified
choice of a majorant.  Finally set
\begin{equation}
 \Theta=E_0-\Omega(V_0+V_1),\qquad
 \boxed{\bigl(\Theta^2-\Omega^2(a_0+a_1)\bigr)^2
                   -4\Omega^4a_0a_1=0.}
 \label{eq:origin-nine}
\end{equation}
The three literal division-free DAGs, all fixed numerical leaves, and
finite coefficient reconstruction recipes are included in the packet.

\begin{theorem}[Preservation of the consistency encoding]
The three positive-domain equations \eqref{eq:origin-nine} have a zero
if and only if their respective inherited positive source systems have
a common zero.  Composing with the retained native-source and
machine/refutation calibration gives the same ordinary theorem
formalizable in ZFC:
\[
 \mathrm{ZFC}\vdash
 \bigl(\nexists(b,g,h,s,w,C,i,j,n)\in\mathbb Z_{>0}^9:
       \mathcal F=0\bigr)\ \longleftrightarrow\ \operatorname{Con}(\mathrm{ZFC}).
\]
\end{theorem}

\begin{proof}
The preceding arguments establish all the numerical prefix hypotheses,
its soundness range, cofinal completeness, and the exact central-binomial
calibration.  The source zero set is unchanged by
Lemma~\ref{lem:radix-origin}.  The saturated positive ProductPell finalizer
is retained unchanged.

For clarity, its proof order matters.  Unconditional strict bounds
$V_0>\sqrt{a_0}$ and $V_1>\sqrt{a_1}$ first force $m>0$ from a norm zero.
They then give the carrier range and strict separation of the two
radicands.  A nonzero rational sum of their signed square roots forces
each root to be rational, hence integral.  This supplies the square
conditions and $J\mid Q$.  The saturated ProductPell theorem recovers
$C=\psi_A(2R+1)$, after which the original numerical ratio theorem
recovers the binary base and \eqref{eq:origin-calibration}.  No inference
about square roots is used prematurely to obtain the margin.
Conversely a source zero supplies fresh ratio and Pell witnesses,
and the inherited integral choice of $n>0$ makes one norm factor vanish.
The recorded soundness and saturation proofs prove these steps in full.

For the simple majorant the identity
\[
 (ADr_0)^2-DF=D(Dr_0^2-1)>0
\]
implies the required strict bound after multiplication by $G^2H_1^2$.
For the complete majorant, the exact polynomial
$\Pi^2-\sigma^2DFI$ has 36,309 strictly positive coefficients after
$A=x+3,C=y+1,i=u+1,Z=v+1$, including a positive constant.
The translated polynomial $\Pi$ itself also has positive coefficients.
Thus the same majorant premise holds uniformly.  These are retained
algebraic certificates, rechecked exactly on the present parameters.

All new steps are elementary finite polynomial identities, integer digit
arguments, or arithmetic statements with the same formalizable inherited
interfaces.  The packet does not supply a new end-to-end Lean/Rocq proof
or an exported object-language ZFC derivation; its calibration dependencies
are explicit rather than inferred from finite testing.
\end{proof}

\subsection{Exact degrees and exact arithmetic-circuit checks}

For the original31 prefix,
\[
 r=\deg R=7931,\quad \nu=\deg N=2106,\quad
 a=\deg A=r+4\nu+2=16357.
\]
For compact54 these values are $8787,2320,18069$.
The positive leading source term is $R\sim 2^{\lambda+1}bB^{2E-1}$,
where $(E,\lambda)=(1983,1)$ or $(2197,0)$.
Thus these bounds are attained on the actual fixed-coefficient instances.
Write $\delta_0=V_0^2-a_0$, $\delta_1=V_1^2-a_1$, and
\[
 B_s=E_0^2-2E_0\Omega(V_0+V_1)+\Omega^2(\delta_0+\delta_1).
\]
The exact cancellation-resolved identity is
\begin{equation}
 \mathcal F=B_s^2+4\Omega^2V_0V_1B_s
                  +4\Omega^4(\delta_0V_1^2+a_0\delta_1).
 \label{eq:origin-degree-identity}
\end{equation}
One has
\[
 \deg\Omega=5a+2W+3,\quad \deg V_0=3a+r+1,\quad
 \deg\delta_0=3a+r+2,\quad \deg E_0=a+4.
\]
For the simple majorant,
$\deg V_1=18a+r+24$ and $\deg\delta_1=34a+2r+48$.
The uniquely dominant term of \eqref{eq:origin-degree-identity} is
$\Omega^4\delta_1^2$, giving
\[
                  \deg\mathcal F=88a+4r+8W+108.
\]
For the complete majorant, exact expansion gives
\[
 \deg V_1=18a+r+30,\qquad \deg\delta_1=18a+2r+60.
\]
Their unique leading monomials before composition are
\[
1048576A^{18}C^{24}i^6Z,\qquad
624687054848A^{18}C^{48}i^{12}Z^2.
\]
The four contribution degrees in \eqref{eq:origin-degree-identity} are
$948768,997810,989850,752426$.  The second is uniquely maximal, and its
leading coefficient is positive, giving degree $997810$.
Merely propagating maximum degrees through the unresolved norm would
not establish these exact values.

The evaluator uses one addition chain for every required power of $B$.
The selected circuits evaluate $\ell$ and $T_*$ as sums of shared powers,
and evaluate the three coefficient kernels in Horner form.  This saves
26 multiplications for the original31 forms compared with using Horner
form for both masks; it saves 25 for the compact profile.  The common
power chain is a finite table of equalities $e=u+v$; its use costs one
multiplication per live chain step.  Fixed numeral construction is excluded
from this convention, but all live binary additions, subtractions,
squares, and multiplications by numerals are counted.

The supplied verifier does not import the builder.  It proves 54 or 55
exact polynomial identities at the named registers of each DAG, including
all complete-majorant coefficients.  It separately checks the variable
set, acyclicity, liveness, operation counts, fixed-leaf recipes, and the
final output.  A second independently implemented sparse ring checks all
22 source identities including absent monomials, all 20,466 original31
and 22,450 compact convolution interactions, the coefficient bounds, the
positive majorant certificate, the population exponents, and the resolved
degree.  Eight complete native coefficient byte strings are reconstructed
from the retained literal finite inputs.  These checks support exact
algebra and the numerical instances; they do not experimentally test
consistency or substitute for the unbounded semantic calibration.

The coefficient-size tradeoff is material.  The original31 tuple occupies
533,099 raw bytes in its four ports, versus 215,492 for compact54.
Its twelve live fixed leaves occupy 4,001,391 raw bytes and 53,739 gzip
bytes, versus eleven leaves occupying 1,565,018 raw bytes and 23,481 gzip
bytes for compact54.  The larger live-leaf total includes the fixed powers
of two and derived coefficient products.  These are distinct measures
from circuit length and from the size of the entire research archive.

\subsection{Further questions suggested by the construction}

The uniform translation is maximal for these fixed supports, but a joint
optimization of coordinate weights, target positions, row operations and
evaluator structure may improve all three circuits further.  Such a search
should allow weight zero from its beginning and score the final polynomial
rather than only the convolution height.  Another concrete target is a
Lean proof of the general layout-translation and zero-position mask lemmas,
composed with the existing ratio interfaces.  Finally, the exact register
certificates suggest a small checked circuit verifier that accepts the
whole finite DAG and the fixed numerical leaves without trusting the
arithmetic-circuit producer.  None of these questions entails a claim that
nine variables are necessary or that a different direct proof encoding
cannot give a simpler equation.

\section{A direct quartic with 495 nonnegative unknowns}
\label{sec:projected-direct}

The direct branch encodes finite derivations in the repository's calibrated
numerical calculus $C_{222}$.  Its immediate parent is the incoming
synchronized-certificate construction with 522 nonnegative unknowns.
Combining its arithmetic certificates with branchwise projection of rule
labels gives the following fully materialized equation.

\begin{center}
\begin{tabular}{lrr}
\hline
Measure & Synchronized parent & Projected equation\\
\hline
Nonnegative unknowns &522&495\\
Quadratic residuals &400&391\\
Exact total degree &4&4\\
Expanded monomials &20,225&15,377\\
Maximum absolute coefficient &83,900,628&83,867,632\\
\hline
\end{tabular}
\end{center}

The actual polynomial is
\[
 P_{495}(z_0,\ldots,z_{494})=\sum_{j=1}^{391}R_j(z_0,\ldots,z_{494})^2.
\]
Every coefficient of every quadratic $R_j$, and every collected coefficient
of $P_{495}$, is supplied as an ordinary integer in the literal artifacts.
The notation does not quantify additional registers.  The source equations
below explain the construction; all their powers, binomials, factorials,
and bounded quantifiers have been removed from the delivered quartic.

This is a measured improvement of the direct, modest-coefficient quartic
branch.  It is not a claim that 495 is an optimal variable count.  The
construction does not encode a run of K215 and does not use Jones's
universal polynomial.  Its ordinary consistency theorem imports the
identified $C_{222}$ calibration; it is not a new end-to-end proof-assistant
certificate.

\subsection{Project the premise and conclusion labels first}

Write
\[
 I(a,b)=2(a+b)(a+b+1)+4b,\qquad A(a,b)=I(a,b)+2.
\]
The calibrated source has twelve branches.  Its truth and axiom branches
have premise labels $p=q=0$ and conclusion label
$v=h_b(a,b,c,d)$, with truth represented by $h_b=0$.
The exact nine nonzero conclusion polynomials are the inherited numerical
constructors.  The two inference branches admit the parameterizations
\[
 (p,q,v)=(a,I(a,b),b),\qquad
 (p,q,v)=(b,0,A(a,b)).
\]
For modus ponens these parameters are the old first premise and conclusion;
for generalization they are its generalization parameter and premise.
Every old instance therefore has a new instance, and substitution recovers
an old instance from every new one.  Since existential quantification
distributes over a finite disjunction, unused parameters in an unselected
branch impose no compatibility condition.

There is one further local saving.  The bundle branch's conclusion is
independent of its fourth parameter.  Its guard
\[
 g^2=(d_{\rm old}+1)^2,
 \qquad g=(c-a)(c-b),\qquad d_{\rm old}\ge0,
\]
can be replaced by $g^2=d+1$, $d\ge0$.  Both existential conditions mean
$g\ne0$.  An old witness gives $d=d_{\rm old}^2+2d_{\rm old}$; a new
witness gives $d_{\rm old}=|g|-1$.  The fourth parameter in another
selected branch retains its original meaning.

Fix outer parameters $t,m\ge1$ and $x,y,M_0\ge0$.  Set $M=M_0+1$ and
$w(j)=1+Mj$.  For each branch let
\[
 e_b=S' w(v_b)-S w(p_b)w(q_b).
\]
Let $e_0,\ldots,e_8$ be the truth and eight unguarded axiom residuals;
let $e_{\rm MP},e_{\rm GEN},e_{\rm B}$ be the two inference and bundle
residuals.  Put $G=((c-a)(c-b))^2-d-1$.  The complete row equation is
\begin{align*}
0={}&\left(e_0\cdots e_8e_{\rm MP}e_{\rm GEN}\right)^2
                   (e_{\rm B}^2+G^2)\\
 &+\bigl[x-S-(1+(i+1)y)u\bigr]^2
  +\bigl[x-S'-(1+(i+2)y)v\bigr]^2.
\end{align*}
Every summand is nonnegative.  Consequently this equation is equivalent
to the two beta reconstructions together with one valid selected rule's
accumulator identity.  Only eight bounded row coordinates remain:
\[
 (a,b,c,d,S,S',u,v),\qquad
 (a,b,c,d,S,S',u,v)\le(m,m,m,m,y,y,x,x)
\]
coordinatewise.  The implementation also shares $S'M$ in
$S'w(h)=S'+(S'M)h$ and squares the product of the eleven unguarded
residuals once.  These are exact polynomial identities, not probabilistic
fingerprint equalities.

\subsection{A small circuit bounds all labels}

At $a=b=c=d=m$, a literal certificate bounds every constructor by a fixed
power of $R_0=m+2$.  Assign exponent one to $m$, exponent
$\lceil\log_2 c\rceil$ to a positive numeral $c$, exponent zero to zero,
and propagate
\[
 e_{U+V}=\max(e_U,e_V)+1,\qquad e_{UV}=e_U+e_V.
\]
Since $R_0\ge2$, induction proves $U\le R_0^{e_U}$ at every gate.
All constructor coefficients are nonnegative, so the diagonal evaluation
also bounds the labels when the four parameters are merely at most $m$.
The largest certified exponent is $3,445,903$; use
\[
 h(m)=(m+2)^{4,194,304},\qquad M=(h(m)+1)^t.
\]
The fixed power defining $h$ takes 22 squarings.  This high semantic degree
does not raise the quartic degree of its arithmetic graph.

Here is the exact carrier argument.  For the $t$ selected rule instances
define
\[
 B_0(X)=\prod_{i<t}(1+v_iX),\qquad
 B_1(X)=(1+X)\prod_{i<t}(1+p_iX)(1+q_iX).
\]
The endpoints $S_0=M+1$, $S_t=1$ and the accumulator identities imply
$B_0(M)=B_1(M)$.  All states are positive, by induction from the first
endpoint and positivity of the label weights.  With $h=h(m)$,
\[
 B_0(M)\le(1+hM)^t<((h+1)M)^t=M^{t+1}.
\]
Thus $\deg B_1\le t$, since a nonnegative integer polynomial of greater
degree would have value at least $M^{t+1}$.  Both polynomials now have
coefficient sum at most $(h+1)^t=M$, and constant coefficient one.  Each
coefficient is therefore strictly smaller than $M$.  Uniqueness of
base-$M$ digits proves $B_0=B_1$.  Comparing the multiplicities of the
distinct factors $1+jX$, $j>0$, gives the exact Horn balance for target
label one.  The inherited closure-separation argument turns that balance
into a finite derivation, without excluding cycles or repeated premises.

Conversely, expand a finite derivation to a finite tree.  Choose $t$ for
its rule occurrences and choose $m$ afterward to bound their four
parameters, including each new bundle slack $g^2-1$.  Products of the
weights of outstanding obligations give the required integral
accumulator trace.  The independent parameter $m$ avoids any need to pad
the derivation until its horizon bounds its own labels.

\subsection{Compile and contract}

Choose $y$ above the finitely many states and divisible by $t!$, then
choose a beta code $x$ for the trace.  Its natural quotients are at most
$x$.  Conversely, the bounds $S,S'\le y$ make both reconstructions
canonical remainders, since $y<1+(i+1)y$.

A second literal majorant certificate uses base $2+t+x+y+m+M_0$ and
fixed exponent $8,388,608$.  It bounds both nonnegative sides of the
row equation at all upper row bounds, as well as the index and the
individual bounds.  The inherited fixed-vector CRT theorem then applies
to the eight columns.  Its factorial factors cancel because the CRT
moduli are one modulo the chosen common factorial.

The resulting finite arithmetic circuit makes 11 binomial requests and
24 positive-power requests, with 19 distinct Pell indices and 10 distinct
power bases.  The shared binomial radix and common Pell-frame theorems
apply to this finite family unchanged.  In soundness one identifies the
power graphs first, the binomial digits second, and the factorial/CRT
semantics afterward.  In completeness one chooses the intended acyclic
semantic values first and the shared auxiliary certificates afterward.
This order avoids circular justification of the shared radix or frame.

The emitted parent of the final quartic has 823 natural variables and
719 quadratic residuals.  Exact nonnegative monic substitutions eliminate
328 further variables and residuals.  The inverse substitutions have
nonnegative integer coefficients, and the retained residuals remain
quadratic.  Thus both existence implications hold over the naturals.
The independent replay verifies all 328 substitutions and the complete
15,377-term sum-of-squares expansion.

\section{Recover nonnegativity from the equations themselves}
\label{sec:frame-domain}

Converting every natural coordinate independently would waste many
variables.  The displayed equations already force most coordinates to
have the required signs.  A finite certificate shows that guarding only
170 of the 495 coordinates suffices.  The main additional argument uses
the common Pell frame before invoking any exponentiation semantics.

\subsection{Local sign implications}

Suppose a literal residual can be written
\[
 A(\mathbf z)z_i=B(\mathbf z),
\]
where neither $A$ nor $B$ involves $z_i$.  If previously established facts
prove $A>0$ and $B\ge0$, then $z_i\ge0$; if $B>0$, then $z_i>0$.
The certificate uses this elementary implication and the weakening
$z_i>0\Rightarrow z_i\ge0$.  Its sign checks allow even powers without
assuming a sign for their bases.  Positive integer monomials can also
pay an explicitly bounded negative constant term.

Starting with the 170 declared guards, 531 such checked implications
prove 447 of the 495 coordinates nonnegative.  Crucially, this first
stage proves all prerequisites of the following common-frame argument.
No $T$ or Pell abscissa whose sign is recovered by that argument is used
prematurely.

\subsection{One frame proves 29 further signs}

The exact raw compiler identities restore polynomially from the final
495 coordinates.  Let $W,A,Z$ be its common frame, and let
$(n,k,m)$ range over the finite power requests.  The frame identities are
\[
 W=2+\sum_{\xi\in\mathcal A}\xi,\qquad
 A^2=1+W^3(W+2)Z^2.
\]
Here $\mathcal A$ is the finite set of distinct base, exponent, and output
atoms in the circuit; distinct atoms with equal values are still counted
separately.  The exact list is supplied.  Before the new argument,
all these atoms are nonnegative, every base $n$ is strictly positive,
$A\ge0$, and $Z\ge1$.  These facts are checked through the literal
restorations, including the factorial output used as a power base.
Consequently
\[
 W\ge2,\qquad 1\le n\le W,\qquad A>W^2.
\]
The last strict inequality follows immediately by squaring; no Pell
classification or already-established power semantics is needed.

Each of the ten distinct base moduli satisfies
\[
 T_n=2An-n^2-1.
\]
Therefore $T_n\ge2A-W^2-1>0$.  For each power request the exact oriented
congruence is
\[
 X+Yn=YA+m+T_nq,
\]
and the first sign stage has already established $Y,m,q\ge0$.
Since $A-n>0$, rearranging gives
\[
 X=Y(A-n)+m+T_nq\ge0.
\]
This recovers the signs of all 19 distinct index abscissae.  After adding
the ten modulus signs and these 19 abscissa signs, 21 further local
implications prove all 495 coordinates nonnegative.

The order matters: base positivity precedes the frame argument, modulus
positivity precedes abscissa recovery, and full natural-domain arithmetic
semantics is invoked only after every coordinate has been recovered.
This separates a short ordered-ring proof from the much stronger
Pell-index argument used later to identify actual powers.

\subsection{A quartic in 977 unrestricted integers}

For an integer $x$, the classical three-square theorem gives
\[
 x\ge0\quad\Longleftrightarrow\quad
 \exists u,v,w\in\mathbb Z:\quad x=u^2+v^2+w(w+1).
\]
Soundness is immediate.  For completeness, $4x+1$ has a representation
by three squares; modulo four exactly one square is odd.  Writing that
odd integer as $2w+1$ and the even integers as $2u,2v$ proves the formula.

Adjoin this quadratic identity for each of the 170 guarded coordinates.
The other coordinates require no separate representation, because the
checked domain-recovery theorem forces their nonnegativity.  Keeping the
original coordinates at this stage preserves degree two in every
residual.  The intermediate unrestricted-integer system therefore has
\[
 495+3\cdot170=1005\quad\text{variables},\qquad
 391+170=561\quad\text{quadratic residuals}.
\]
Twenty-eight further signed monic substitutions preserve quadratic
degree and existence over the integers.  The final literal equation has
\[
 \boxed{\begin{gathered}
 977\ \text{unrestricted integer unknowns},\qquad
 533\ \text{quadratic residuals},\qquad \deg P_{977}=4,\\
 24,632\ \text{expanded monomials},\qquad
 \max|\text{coefficient}|=83,869,250.
 \end{gathered}}
\]
The comparison with the synchronized parent's 2088-integer quartic is
a reduction of 1111 integer variables, or approximately $53.2\%$, while
retaining degree four.  This is a comparison within the direct branch;
the repository's much smaller-variable high-degree constructions remain
different tradeoffs.

Every integer zero of $P_{977}$ restores the guarded 495-coordinate
system, hence a completely nonnegative zero of $P_{495}$ and a calibrated
finite refutation.  Conversely, a refutation gives a natural zero,
three-square witnesses for the 170 guards, and then a zero of the
contracted integer system.  Composing with the named ordinary source
calibration and negating yields the desired ZFC-provable biconditional
between absence of integer zeros and $\operatorname{Con}(\mathrm{ZFC})$.

\subsection{What the finite verifier establishes}

The domain checker is independent of the sign-discovery code.  It
rederives every local implication from a literal quadratic residual and
checks 1535 raw polynomial restorations.  Of the raw equations, 1040
restore to zero identities and 391 restore to final residuals.  It also
binds all 1243 raw arithmetic-gate metadata records to their exact
defining residuals before expanding them.  Finally it checks the frame
norm and sum, ten modulus identities, and all 24 oriented power
identities.  The ordinary inequalities above then justify its 29
additional sign conclusions.

Corruption controls require rejection after a gate operand is changed,
a guard is deleted, a pivot orientation is reversed, a frame conclusion
is omitted, or a modulus identity is replaced.  Normal and optimized
Python runs perform the checks explicitly.  These finite verifications
certify the formulas, transformations, and coefficient data; they do not
replace the unbounded source proof or decide whether a refutation exists.


\section{Logical dependencies and reproducibility}
\label{sec:dependencies}

\subsection{The direct branch reaches consistency without a machine run}

The direct source is the numerical calculus $C_{222}$, with distinguished
contradiction code one. Two inherited theorems identify its consistency
meaning~\cite{calibration}. In the retained source, the first theorem is
\begin{quote}\small\ttfamily
KEndpoint.c222\_iff\_zf\_refutation
\end{quote}
and states that derivability of the distinguished code in $C_{222}$ is
equivalent to a finite ZF refutation. It appears before the
machine-parameter section of \texttt{KEndpoint.lean}; it has no machine,
tape, simulation, boot, coverage, or halting premise. Its upper direction
uses the interpretation of the numerical rules, and its lower direction
uses the established semantic completeness chain of the calculus.

The second theorem is
\begin{quote}\small\ttfamily
ConstructibleChoice.zf\_refutation\_iff\_zfc\_refutation.
\end{quote}
Its forward implication uses theory inclusion. Its reverse implication
uses the constructible model of a model of ZF, together with first-order
soundness and completeness. These are the inherited ordinary calibration
results needed for the direct polynomial. Their complete immediate
declaration files and their exact source hashes are retained under
\path{direct/provenance/upstream/calibration/}.

The retained calibration revision is
\texttt{08200eb10fde2db408c69dffd840d6920a569b1f}. This is intentionally
distinct from the comparison snapshot in the title: the incoming parent
pinned those earlier theorem bodies, and this packet preserves its pins
instead of silently replacing them with a later branch. The ordinary CRT,
binomial, factorial, power, and shared-frame proofs are also retained.

Thus the direct branch encodes derivations in a numerical calculus and
does not simulate K215 or substitute parameters into Jones's universal
polynomial. Some files in the inherited Lean import graph also support
machine constructions. No claim is made that the entire import graph is
free of machine-related modules.

\subsection{The compact branches preserve their native calibrations}

The fourteen-integer branch retains the coherent source-34/source-71
coefficient tuple and its signed native kernel, ordinary compiler
calibration, and K215/refutation endpoint. The new proof establishes a
two-way existence equivalence to that kernel, including the restoration
of its strong variables. The six fixed coefficient values are finite
literal integers in the packet. Their finite reconstruction is checked;
their semantic interpretation comes from the inherited compiler proof.

The nine-variable branch likewise keeps its two coherent native tuples
separate: the original31 \texttt{gapN86} source and the compact54
\texttt{gapNq85} source. It imports the positive sparse-carry encoding,
the numerical ratio theorem, the saturated ProductPell theorem, and the
native source/refutation calibration. The new proof checks the exact
premises affected by the origin shift and recalibrates the threshold;
it does not infer those unbounded implications from a few sampled values.

The conversion from a nonnegative integer to three integer witnesses is
classical. For $x\in\mathbb Z$, the exact characterization
$(\exists u,v,w\in\mathbb Z:\ x=u^2+v^2+w^2+w)\Longleftrightarrow x\ge0$
is also recorded in the
Archive of Formal Proofs formalization of the three-square theorem
\cite{three-squares}. Here it is used only after the 170-coordinate domain
recovery theorem has been established.

\subsection{Four distinct levels of evidence}

It is useful to separate what has actually been delivered.
\begin{enumerate}
 \item \textbf{Explicit equations.} All seven final polynomials are
 literally determined. The large-coefficient forms have acyclic circuits
 and materialized fixed leaves. Both quartics have complete expanded
 coefficient lists as well as quadratic sum-of-squares presentations.
 \item \textbf{Ordinary mathematical proofs.} The article proves the
 new source changes, guard and radical arguments, domain recovery, and
 witness restorations. The retained sources state the precise inherited
 semantic calibrations used in composition.
 \item \textbf{Exact finite certificate checks.} Independent arithmetic
 implementations check relevant identities, source layouts, leading
 terms, sign implications, contractions, and collected coefficients.
 Rebuilding recovers the delivered equation files byte for byte.
 \item \textbf{Formal proof artifacts.} Existing pinned Lean declarations
 are retained as dependencies. This work has not built a new kernel-checked
 theorem for the seven final polynomials or exported a derivation in the
 sealed object-language ZFC system. Those remain concrete formalization
 tasks, not claims supported by the Python checks.
\end{enumerate}

The internal mathematical reviews in the packet were conducted by separate
AI work streams. They are useful adversarial checks of the construction,
not external peer review. The finite verifiers themselves are short
ordinary programs and remain subject to inspection; their rejection
controls exercise actual failure paths, including when Python assertions
are disabled.

\subsection{Reproduce all finite results}

Python's standard library is sufficient. From the extracted packet root,
run
\begin{Verbatim}[fontsize=\small]
python3 verify_all.py --output verification/rechecked.json
python3 -O verify_all.py --output verification/rechecked_O.json
\end{Verbatim}
The driver copies the three components to a temporary directory, checks
the delivered artifacts before rebuilding them, and compares all
reconstructed final equations with the original bytes. It leaves the
delivered equation and certificate files unchanged. The two recorded root
receipts are \path{verification/ROOT_NORMAL.json} and
\path{verification/ROOT_OPTIMIZED.json}.

The finite checks include the following substantive obligations:
\begin{itemize}
 \item All 495 coefficients of the four-root norm, four exact root-defect
 identities, both degree-leading forms, all six fixed coefficient values,
 and the literal fourteen-variable circuits.
 \item All 22 source extraction identities for each positive packing,
 including absent monomials; 20,466 and 22,450 convolution interactions;
 all eight native coefficient ports; and 54 or 55 exact register identities
 per final nine-variable circuit.
 \item The complete 36,309-term positivity certificate for the full
 radical majorant and the cancellation-resolved degree calculation.
 \item The direct source constructor identities, both majorants, arithmetic
 certificate interfaces, 328 nonnegative contractions, and the full
 quartic expansion.
 \item The 170-guard sign certificate, 1,535 raw coordinate restorations,
 1,243 gate bindings, the common frame and power identities, 28 signed
 integer contractions, and all 24,632 final quartic coefficients.
 \item Rejection after a guard, gate operand, pivot orientation, frame
 conclusion, modulus identity, restoration coefficient, or expanded
 coefficient is corrupted.
\end{itemize}
Finite modular comparisons retained by the fourteen-variable audit are
transcription diagnostics. They are not substituted for exact symbolic
identities, degree arguments, or unbounded semantic proofs.

\section{Repository integration and next mathematical targets}
\label{sec:next}

\subsection{Files to treat as authoritative}

The root \texttt{README.md} maps each construction to its literal file.
For $P_{540}$ and $P_{568}$ use the JSON circuits immediately under
\path{integer14/}; for the three nine-variable forms use
\path{nine/equations/}. The natural quartic is
\path{direct/equations_v3/cap64/equation.json}, and the integer quartic is
\path{direct/integer977/equation.json}. Each quartic directory also contains
\texttt{expanded\_polynomial.json}. The accompanying schema note specifies
coefficient signs, variable indices, fixed leaves, and monomial conventions.

The archive keeps original source notices and the repository's MIT No
Attribution license. The comparison ledger records the inspected main
commit, incoming archive names, Git object identifiers, and downloaded
SHA-256 values. Original archive filenames containing a later date have
not been renamed. The packet-wide SHA-256 manifest binds the final files;
it is an integrity record, not a mathematical proof.

No existing numbered repository construction is silently replaced. A new
entry should state the domain and all selected size measures explicitly.
The word ``simpler'' should be qualified by those measures. In particular,
the 420-operation original31 form uses a larger native coefficient tuple
than the 428-operation compact54 form. Keeping both makes that choice
visible instead of hiding the coefficient cost in fixed leaves.

\subsection{Research questions made concrete by this packet}

\paragraph{Can the four-root mechanism keep thirteen variables?}
The gain at fourteen variables comes from retaining the guard root.
Reabsorbing it without reintroducing the removed amplifier is a precise
target. Any proposed fifth-root fusion must handle dependent square
classes and both branches $cT-R=0$ and $cT-R\ne0$; a generic norm identity
alone will not establish rational extraction.

\paragraph{Can a better center reduce degree below 540?}
The exact defect identities explain which factors remain expensive.
They provide a starting point for a systematic search over polynomial
root approximations. A candidate center must retain the CRT congruences,
integrality of the final witness, and control of exceptional branches.
Degree estimates must be resolved after cancellation on the authentic
fixed-coefficient slice.

\paragraph{Can source layout and evaluator be optimized jointly?}
The uniform radix shift is maximal in its restricted class. That leaves
nonuniform weight changes, target changes, source row transformations,
and globally shared addition chains. A useful search objective is the
actual final circuit and degree, with coefficient bytes reported alongside
them. Optimizing convolution height alone can miss an evaluator saving.

\paragraph{How few signs really need guards?}
The set of 170 guards is sufficient, not proved minimal. A smaller seed
set for the present sign rules, additional multi-row inequalities, or
stronger frame consequences may save more variables. The explicit
certificate separates discovery from verification, so a heuristic search
may propose guards while an independent checker establishes the result.
The final count should be measured after signed contractions, since a
guard's cost can depend on the resulting elimination opportunities.

\paragraph{Can the direct source be reduced before arithmetic compilation?}
The successful projection removes deterministic labels before applying
CRT. Further source changes should preserve exact parameter dependencies,
the finite-tree completeness argument, and the carrier coefficient bound.
Sharing expressions across mutually exclusive rule branches may give
additional savings without changing the underlying numerical calculus.

\paragraph{Can the whole endpoint be checked in one proof kernel?}
A particularly valuable next deliverable would connect the emitted
literal quartic, its coefficient expansion, the domain certificate, and
the inherited $C_{222}$ consistency endpoint in a single kernel-checked
theorem. The certificate formats and pinned declarations in this packet
make that task specific. It requires formalizing the universal arguments,
not merely embedding successful Python exit codes in a theorem statement.

\begin{thebibliography}{9}
\raggedright

\bibitem{proveit}
Vladimir Reshetnikov and ProveIt contributors.
\emph{ProveIt: Hilbert's tenth problem, Diophantine consistency forms,
and incoming proposals.}
Comparison commit \texttt{7e533c33819ee1f703b69f2778a767cdc069940d}.
\url{https://github.com/VladimirReshetnikov/ProveIt/tree/7e533c33819ee1f703b69f2778a767cdc069940d/docs/incoming}.
The accompanying comparison ledger records individual archive hashes.

\bibitem{calibration}
ProveIt contributors.
\emph{KEndpoint and ConstructibleChoice calibration sources.}
Retained revision \texttt{08200eb10fde2db408c69dffd840d6920a569b1f}.
\url{https://github.com/VladimirReshetnikov/ProveIt/blob/08200eb10fde2db408c69dffd840d6920a569b1f/Computability/BusyBeaver/Lean/BusyBeaver/ZFC/KEndpoint.lean}.
Exact immediate copies and SHA-256 pins are included in the packet.

\bibitem{three-squares}
Anna Danilkin and Lo\"ic Chevalier.
\emph{Three Squares Theorem.}
Archive of Formal Proofs, 3 May 2023.
\url{https://isa-afp.org/entries/Three_Squares.html}.

\bibitem{bayer}
Jonas Bayer, Marco David, Malte Ha\ss ler, Yuri Matiyasevich, and
Dierk Schleicher.
\emph{Diophantine equations over $\mathbb Z$: universal bounds and
parallel formalization.}
arXiv:2506.20909v3, 22 October 2026.
\url{https://arxiv.org/html/2506.20909v3}.

\end{thebibliography}
\end{document}
